% Benötigt: \usepackage{booktabs,array}
% Stand: Repo AlexanderGiller/master-thesis-proover (main). Zeilenzahlen mit `wc -l` vor der Abgabe aktualisieren.
% Mit "CODE:" markierte Stellen weichen vom Design-Kapitel ab und müssen vor der Evaluation angepasst werden.

\section{Implementation}
\label{sec:implementation}

This chapter describes how the design of Chapter~\ref{sec:design} is realized in Python.
Section~\ref{subsec:impl-overview} gives an overview of the code base.
Sections~\ref{subsec:impl-model} and~\ref{subsec:impl-parser} describe the data model and the parser, Section~\ref{subsec:impl-driver} the driver that dispatches the steps and computes the verdict.
Sections~\ref{subsec:impl-global} to~\ref{subsec:impl-atp} cover the checkers, Section~\ref{subsec:impl-resources} the resource control, and Sections~\ref{subsec:impl-testing} and~\ref{subsec:impl-tooling} the test suite and the tooling used for the evaluation.
Section~\ref{subsec:impl-limitations} lists where the implementation falls short of the design.

\subsection{Overview}
\label{subsec:impl-overview}

The implementation is written in Python~3.12 and uses the parser library \texttt{lark} as its only runtime dependency.
It consists of about 5\,000 lines of code in four layers (Table~\ref{tab:impl-modules}): the parser and the AST, the rule checkers with shared utilities, the driver, and the tooling for scoring and performance measurement.
The separation follows the responsibilities of the design: syntax handling is independent of logical verification, and each rule family lives in its own module, so that new constructs or rules can be added without touching the others.

\begin{table}[t]
\centering
\small
\caption{Module structure of the implementation (approximate line counts).}
\label{tab:impl-modules}
\begin{tabular}{@{}l p{0.60\linewidth} r@{}}
\toprule
Layer & Module (\texttt{src/\dots}) & Lines \\
\midrule
Parser and AST & \texttt{parser/parser.py}, \texttt{parser/tptp\_grammar.lark}, \texttt{parser/ast\_nodes.py} & 560 \\
Constants & \texttt{var\_mapping.py} (enumerations for roles, statuses, rules, connectives) & 90 \\
Global checks & \texttt{checker/file\_dir.py}, \texttt{dependency\_checker.py}, \texttt{negated\_conjecture\_checker.py} & 500 \\
Structural and propositional & \texttt{structural\_checker.py}, \texttt{propositional\_checker.py}, \texttt{junction\_checker.py} & 430 \\
Quantifier rules & \texttt{instantiate\_checker.py}, \texttt{existential\_gen\_checker.py}, \texttt{skolem\_checker.py} & 760 \\
Clausal and equational & \texttt{clausal\_checker.py} & 470 \\
Shared utilities & \texttt{alpha\_eq.py}, \texttt{formula\_utils.py} & 400 \\
ATP fallback & \texttt{external\_atp\_checker.py} & 190 \\
Driver & \texttt{main.py}, \texttt{deep\_recursion.py} & 1\,130 \\
Tooling & \texttt{scoring/}, \texttt{performance/} & 650 \\
\bottomrule
\end{tabular}
\end{table}

% TODO: Abbildung "Modulabhängigkeiten" im Ordner diagrams zeichnen (parser -> ast_nodes; checker/* -> ast_nodes, alpha_eq, formula_utils; main -> parser, checker/*, external_atp_checker).

\subsection{Internal Representation}
\label{subsec:impl-model}

Formulas are represented by Python data classes defined in \texttt{ast\_nodes.py}.
Terms are \texttt{Variable}, \texttt{Constant}, and \texttt{FunctionTerm}; atomic formulas are \texttt{Atom} and \texttt{Equality}; compound formulas are \texttt{Negation}, \texttt{BinaryFormula}, \texttt{JunctionFormula}, and \texttt{QuantifiedFormula}.
Conjunctions and disjunctions are \texttt{JunctionFormula} objects with a list of operands rather than nested binary trees, which simplifies the junction and clausal checks, since many rules operate on lists of literals.
A \texttt{QuantifiedFormula} stores the quantifier and the list of bound variables, which the instantiation, generalization, and Skolemization checks rely on.

A proof step is an \texttt{AnnotatedFormula} with name, role, formula, and either a source or an \texttt{InferenceRecord}.
The record contains the rule name, the parents, and a list of information items: the status, the \texttt{new\_symbols} declaration (kind and symbols), and, for Skolemization, the \texttt{skolemize} information with the eliminated variable and the Skolem term.
Roles, statuses, rule names, quantifiers, and connectives are \texttt{StrEnum} classes in \texttt{var\_mapping.py}.
This avoids inconsistencies caused by ad hoc string handling and makes the set of recognized values explicit.

\subsection{Parser}
\label{subsec:impl-parser}

The parser is generated by Lark from the grammar \texttt{tptp\_grammar.lark} using the LALR algorithm with cached tables.
The parser object is created once per process through a cached helper.
A \texttt{TPTPTransformer} maps the grammar productions to the data classes above.
During this mapping some constructs are normalized: for instance, the connective \verb|<=| is converted into a canonical implication direction, so that later checks have fewer cases.
The parser also handles source annotations, inference records, and \texttt{include} directives.
The function \texttt{load\_proof} combines parsing with the extraction of the referenced problem and returns a \texttt{ProofFile} object.

Because the grammar is recursive, deeply nested formulas lead to deep recursion in the transformer and in all checkers.
Section~\ref{subsec:impl-resources} describes how this is handled.

\subsection{Driver and Dispatch}
\label{subsec:impl-driver}

The driver is the function \texttt{check\_proof\_file} in \texttt{main.py}.
It parses the problem file and the proof, builds an index from step names to steps, runs the global checks, and then iterates over the steps in file order.
For each step it selects a checker by the step's role and rule name in an \texttt{if}/\texttt{elif} chain: axioms and conjectures go to the provenance check, the negated conjecture to its dedicated checker, the rules of Table~\ref{tab:rule-families} to their modules, \texttt{skolemize} to the Skolem checker, and every other step with status \texttt{thm} or \texttt{cth} to the external ATP.
Steps with another status and no checker, and steps without any justification, are reported as having no checker.
% CODE: Die if/elif-Kette (ca. 500 Zeilen) soll durch die Registry aus dem Design-Kapitel ersetzt werden (Regelname -> Checker, Arity, geforderter Status).

The driver does not stop at the first error.
Each check writes its findings to a result map keyed by step name; there is one map per rule family (for example \texttt{skolem\_issues} or \texttt{resolution\_issues}), one for the dependency checks, and one for the ATP fallback.
The maps are returned to the caller, so that the console report, the summary table, and the tests all work on the same data (Section~\ref{subsec:design-diagnostics}).
The driver keeps two sets for the freshness check of Skolem symbols: the function symbols of the problem file, collected once, and the Skolem symbols introduced so far.
A \texttt{skolemize} step is rejected if its new symbol is in one of these sets.

\paragraph{Verdict.}
The checker \texttt{check\_proof\_file} returns the issue maps together with the flag \texttt{all\_ok}, which is true if all maps are empty.
The reporting function maps this flag to the SZS status \texttt{VerifiedGood} or \texttt{VerifiedBad}; a timeout yields \texttt{Timeout}.
% CODE: Dreiwertiges Verdikt umsetzen. Aktuell (a) landet ein ATP-Ergebnis "unknown" (ATP fehlt / Timeout / kein SZS-Status) in external_atp_issues und führt zu VerifiedBad,
%       (b) führen Ausnahmen im Worker zu VerifiedBad, (c) gibt es keine eigene Kategorie "unresolved". Gewünscht: unresolved -> uncertain (Ausgabe wie Timeout), nur konkrete Defekte -> failed.
%       Danach diesen Absatz auf die dreiwertige Abbildung aus Tabelle~\ref{tab:verdicts} umschreiben.

\subsection{Global Checks}
\label{subsec:impl-global}

The provenance check in \texttt{file\_dir.py} verifies that an imported axiom or conjecture carries a \texttt{file(\dots)} annotation, that the cited file and formula name exist in the problem, that the role matches, and that the formula is alpha-equivalent to the original.
Alpha-equivalence is implemented in \texttt{alpha\_eq.py} as a simultaneous traversal with two mappings between bound variables.
Before the traversal, consecutive quantifiers of the same kind are merged into one prefix and nested junctions of the same connective are flattened (Figure~\ref{fig:impl-alpha}).
The comparison is used for provenance, for the structural rules, and for many quantifier rules.

\begin{figure}[t]
\centering
\begin{verbatim}
def is_alpha_equivalent(f1, f2):
    return _alpha_eq(normalize(f1), normalize(f2), {}, {})

def _alpha_eq(a, b, map1, map2):
    if type(a) is not type(b): return False
    if a is a Variable:                       # bound or free
        if a in map1 or b in map2:            # bound: mappings must agree
            return map1.get(a) == b and map2.get(b) == a
        return a == b                         # free: names must be equal
    if a is a Quantified formula:
        require same quantifier and arity
        extend map1, map2 with the bound variables of a and b
        result = _alpha_eq(a.body, b.body, map1, map2)
        restore map1, map2                    # scoping
        return result
    otherwise: compare symbols, then recurse over the children in order
\end{verbatim}
\caption{Structure of the alpha-equivalence check (pseudocode, simplified).}
\label{fig:impl-alpha}
\end{figure}

The dependency checks in \texttt{dependency\_checker.py} detect cycles with an iterative depth-first search with three colors and report the cycle; dangling parents are ignored there and reported by the step checks.
A second function rejects any step other than the negation step that cites a \texttt{conjecture}.
The negated conjecture checker verifies role, status \texttt{cth}, the parent, and that the formula is the negation of the conjecture, up to renaming of bound variables.

\subsection{Structural, Propositional, and Quantifier Rules}
\label{subsec:impl-simple-rules}

The structural rules (\texttt{copy}, \texttt{duplicate}, \texttt{rename\_variable}, \texttt{double\_negation}, \texttt{remove\_double\_negation}) are checked in \texttt{structural\_checker.py}: the checker verifies the metadata and that the conclusion has the expected relation to its parent, mostly by alpha-equivalence.
The propositional rules (\texttt{modus\_ponens}, \texttt{conjunction}, \texttt{split\_conjunct}) and the junction rules (\texttt{weaken}, \texttt{commute}, \texttt{excluded\_middle}) act on the top-level connective structure, where the explicit junction nodes pay off.
The syntactic checks for \texttt{weaken}, \texttt{commute}, and \texttt{excluded\_middle} cover only the standard shapes of the rules.
If such a check fails and the step has status \texttt{thm}, the driver asks the external ATP before it reports a failure, so that valid but unusual variants are not rejected.

The \texttt{instantiate} and \texttt{existential\_gen} checks use template matching (\texttt{match\_template\_to\_instance} in \texttt{formula\_utils.py}): the quantified parent is treated as a template with abstract variables, and the checker searches for a consistent substitution that turns it into the conclusion, or the other way around for generalization.
Shared utilities for flattening quantifier prefixes, collecting function symbols, and applying substitutions are collected in \texttt{formula\_utils.py}, so that all checkers interpret formulas consistently.

\subsection{Skolemization}
\label{subsec:impl-skolem}

The Skolem checker (\texttt{skolem\_checker.py}, the largest checker) implements the conditions of Section~\ref{subsec:design-skolem}.
It first validates the shape of the step: rule \texttt{skolemize}, exactly one parent, exactly one status equal to \texttt{esa}, exactly one \texttt{new\_symbols} declaration of kind \texttt{skolem} with one symbol, and a \texttt{skolemize} information item that names the eliminated variable.
It then locates the quantifier that binds the eliminated variable in the parent, analyzing its syntactic and its effective quantifier kind while tracking the polarity (negations and the antecedent of implications flip it), and computes the universal variables in whose scope it occurs.
Finally it checks that the Skolem term is built from the declared symbol and exactly these variables, and that applying the replacement to the parent yields the conclusion up to alpha-equivalence; for this comparison, bound variables are renamed apart first to avoid capture.
The freshness of the symbol is not decided in the checker itself but in the driver, which holds the symbol sets across steps (Section~\ref{subsec:impl-driver}).

\subsection{Equality and Clausal Reasoning}
\label{subsec:impl-clausal}

The module \texttt{clausal\_checker.py} implements \texttt{resolution}, \texttt{paramodulation}, \texttt{reflexivity}, \texttt{transitivity}, and \texttt{rewrite}.
Resolution decomposes the parents into lists of disjuncts, searches for complementary literals, computes unifiers where necessary, builds the resolvent, and compares it to the stated conclusion.
Quantified clauses are handled by flattening the leading universal prefix and rebuilding it for the conclusion.
Paramodulation rewrites occurrences of one side of an equation from one parent in the other parent and compares the result with the conclusion; bound variables are renamed apart before rewriting to avoid variable capture.
Reflexivity checks for a trivial equality $t = t$, and transitivity checks that two parent equalities chain to the conclusion.

The \texttt{rewrite} rule is the weakest check.
In the benchmark format, rewrite steps may omit the demodulating equation, so the intended transformation cannot be reconstructed fully, and the checker verifies only the limited behavior observed in the benchmarks.
% CODE/TODO: Entscheiden, ob rewrite (a) weiterhin nur das beobachtete Verhalten prüft und das in der Discussion als Soundness-Lücke genannt wird, oder (b) bei nicht rekonstruierbarem Schritt an den ATP delegiert bzw. unresolved liefert (Letzteres ist konsistent mit der fail-closed-Politik).

\subsection{Delegation to the External ATP}
\label{subsec:impl-atp}

The module \texttt{external\_atp\_checker.py} implements the fallback of Section~\ref{subsec:design-atp}.
The function \texttt{status\_allows\_atp} implements the eligibility rule: only the statuses \texttt{thm} and \texttt{cth} may be delegated.
The function \texttt{validate\_step\_with\_atp} serializes the premises as axioms and the step as conjecture into a TPTP problem (with quoting of symbols and variable names that are not valid TPTP tokens), writes it to a temporary file, and starts the prover as a subprocess with a time limit.
The default prover is Vampire in CASC mode; the command can be configured through the environment variables \texttt{EXTERNAL\_ATP\_COMMAND} and \texttt{EXTERNAL\_ATP\_BINARY}, and the module translates file paths if the prover is launched through WSL.
The answer is read from the SZS status line of the output (Figure~\ref{fig:impl-atp}).

\begin{figure}[t]
\centering
\begin{verbatim}
validate_step_with_atp(step, premises, timeout) -> True | False | "unknown"
    write premises as axioms, step as conjecture to a temporary .p file
    run prover (subprocess, timeout)
    executable not found        -> "unknown"
    timeout                     -> "unknown"
    no "% SZS status" line      -> "unknown"
    Theorem | Unsatisfiable | ContradictoryAxioms -> True
    any other status            -> False
\end{verbatim}
\caption{Result mapping of the ATP fallback.}
\label{fig:impl-atp}
\end{figure}

The function returns a tri-state value, so that the caller can distinguish a refutation from an unresolved call.
The prover receives only the parents of the step, never the original conjecture or the rest of the proof.
The calls are sequential, and each call may use the full per-proof time limit.
% CODE: Rest-Budget statt festem PER_PROOF_TIMEOUT_SECONDS pro Aufruf (siehe Design, Budget).
% CODE (wichtig): In main.py gibt es einen "shortcut_enabled"-Pfad für consequence/add_hypothesis, der an Dateinamen (PRV013+1.s, PRV015+1.s) bzw. den Pfad /examples/correct/ gekoppelt ist,
%   und die Bedingung "isinstance(parents[0].formula, Negation) != isinstance(parents[1].formula, Negation)" akzeptiert $false auch ohne komplementäre Prämissen. Das ist weder sound noch begründbar;
%   entfernen oder auf den komplementären Fall (_formulas_are_complements) beschränken und nicht dateinamenabhängig machen. Im Text NICHT als Feature beschreiben.

\subsection{Resource Control}
\label{subsec:impl-resources}

Each proof is checked in a separate process created with Python's \texttt{multiprocessing} module.
The parent waits for the result for the configured number of seconds (30 by default, configurable on the command line) and terminates the worker if the limit is exceeded.
The worker passes its result back through a queue; exceptions are caught and returned as an error message.
Process isolation contains crashes and memory exhaustion of a single proof and allows non-cooperative code, such as a hanging ATP call, to be terminated reliably.

Inside the worker, the checking runs in a dedicated thread with a stack of 64\,MB and a recursion limit of 100\,000 (\texttt{deep\_recursion.py}).
This is necessary because both the parser's transformer and the checkers traverse formulas recursively; the benchmark contains stress tests with hundreds of nested negations or function applications.
Raising only the recursion limit would risk a hard interpreter crash, because each Python frame also needs space on the operating system's thread stack.
The dependency analysis uses an iterative search and is therefore independent of proof length.

\subsection{Test Suite}
\label{subsec:impl-testing}

The test suite consists of about 250 \texttt{pytest} test functions in 21 files.
Most files exercise one module in isolation, for example \texttt{test\_skolem\_checker} or \texttt{test\_clausal\_checker}, which is possible because the checkers are functions of the step and its parents.
In addition, there are tests for the parser and the AST, a parametrized provenance test over the example proofs (\texttt{test\_all\_examples}), tests for the tri-state ATP result (\texttt{test\_external\_atp\_checker}, \texttt{test\_main\_atp\_tristate}), edge-case tests, and a mutation-based test (\texttt{test\_mutation\_fuzzing}) that perturbs seed proofs and checks that the verifier reacts with the expected outcome or issue category.
A GitHub Actions workflow runs the suite on Python~3.12 for every push and pull request and fails if the line coverage is below 85\,\% (NFR7).
% TODO: Anzahl der Tests und Coverage aus dem aktuellen CI-Lauf eintragen (pytest --collect-only -q; coverage report).

\subsection{Tooling for Reporting and Evaluation}
\label{subsec:impl-tooling}

The command-line interface in \texttt{main.py} checks either a single proof or, by default, all proofs of the \texttt{ProoVer2026} directory.
For each proof it prints a report grouped by check and finishes with an SZS line, and for a batch it prints counts of the three outcomes and a per-proof summary table with the failed checks.
The scoring module (\texttt{src/scoring}) compares the outcomes with expected results given in CSV files and computes the points of the ProoVer scheme ($+2$, $+1$, $0$, $-1$, $-10$) per proof and in total, together with a breakdown of errors by check category.
The performance module (\texttt{src/performance}) records wall-clock and CPU time per operation category and writes them to a CSV file.
These two modules produce the numbers reported in Chapter~\ref{sec:evaluation}.

\subsection{Deviations from the Design}
\label{subsec:impl-limitations}

The implementation does not yet realize all aspects of the design:
the dispatcher is an \texttt{if}/\texttt{elif} chain instead of a registry;
the verdict is derived from a Boolean flag instead of the three values of Section~\ref{subsec:design-verdict}, so that unresolved ATP calls are reported like failures;
the global check that the derivation ends in \texttt{\$false} (FR9) is not implemented;
\texttt{new\_symbols} annotations of kind \texttt{herbrand} are parsed but not checked separately;
and the \texttt{rewrite} rule is checked only for the shapes observed in the benchmark.
% TODO: Diesen Abschnitt streichen bzw. auf "Known limitations" kürzen, sobald die Punkte im Code behoben sind; die Auswirkungen gehören in die Discussion.